\scanpage{148}
\mspage{64}

ce qui définit bien un foncteur

\[
M^{\mathrm{o}}\times N^{\mathrm{o}}\longrightarrow\mathrm{Ens}.
\]

Malheureusement, pour un des arguments \inserted{fixé}, disons pour $y$ fixé, \deleted{\unsure{c'est}} ce foncteur \deleted{\unsure{un foncteur}}\inserted{n'est} \uline{pas} \deleted{\unsure{limite}} compatible aux limites projectives, p.ex. il ne transforme pas l'objet final de $N^{\mathrm{o}}$ (correspondant à l'objet initial de $N$) en objet final de $\mathrm{Ens}$, mais en \deleted{$\unsure{\operatorname{Hom}(}$} $t(y)=\operatorname{Hom}_{N}(y,t)$. Néanmoins, je pense que

% slanted
\margin{également faux, cf. plus bas}

$s\boxtimes t$ doit bel et bien se définir comme produit de deux objets de $M\boxtimes N$, dépendant fonctoriellement de \unsure{$s$ resp. de $t$|$s$ resp de $t$}. On est donc \unsure{amené|ramené} dès à présent, qu'on le veuille ou non, \deleted{\unsure{à ce que}}, \unsure{à définir|à la question|à la définition} des foncteurs canoniques, commutant aux petites $\unsure{\underrightarrow{\lim}|\underline{\lim}}$

\[
M\longrightarrow M\boxtimes N\longleftarrow N,
\tag*{(4.2.\unsure{14|13})}
\]

% slanted
\margin{\deleted{\unsure{Le}} Cas où $M,N$ sont des catégories-topos nous montre que ces foncteurs doivent être décrits \unsure{comme} $x\longmapsto x\boxtimes e_{N}$, et $y\longmapsto e_{M}\boxtimes y$. Mais il n'est pas clair \unsure{qu'on ait|qu'il y ait}, en dehors du cas des topos, et en admettant l'existence des $x\boxtimes y$, la formule (valable pour les topos) $x\boxtimes y\simeq(x\boxtimes e_{N})\times(e_{M}\boxtimes y)$}

(lesquels seront les foncteurs images inverses des projections canoniques

\diagram{\begin{tikzcd}[ampersand replacement=\&,column sep=large,row sep=small]
\&\unsure{X|\mathcal X}\\
X\times Y\arrow[ur]\arrow[dr]\&\\
\& Y
\end{tikzcd}}

).

4.3. \uline{L'opération} $M\boxtimes N$ \uline{et l'axiome d'accessibilité}.

J'ai passé la journée hier, essayant d'établir l'existence des $x\boxtimes y$ dans $M\boxtimes N$, sans y parvenir -- du moins pas sans faire \deleted{\unsure{des}} une hypothèse supplémentaire d'accessibilité,
% archive sheet number: 147
